\subsection{Constraints of Element Ingestion and Runtime Engines}

A set accepts elements only under specific rules. These rules are not optional decorations around the syntax. They are required by the internal table engine that makes set membership efficient.

When a value is inserted into a set, Python must decide whether that value is already present. It does this through the hash and equality contract introduced earlier. The hash value guides Python toward a candidate region of the internal table. Equality then confirms whether an existing element should be treated as the same logical element.

Therefore, set ingestion has two central constraints:

\begin{itemize}
    \item the element must be hashable;
    \item equal elements must collapse into one stored set member.
\end{itemize}

This is why a set is not just a list without duplicates. A list can accept almost any object reference and preserve repeated occurrences. A set accepts only hashable elements and organizes them through a hash-table engine.

\subsubsection{The Uniqueness Invariant: Automated De-Duplication Mechanics at Runtime}

A set stores each equal value only once. If the same value is supplied several times, Python does not store several independent copies of that value as separate set elements.

\begin{verbatim}
words = {"spam", "eggs", "spam", "spam", "larch"}

print(f"words: {words!r}")
print(f"type(words): {type(words)}")
print(f"len(words): {len(words)}")
\end{verbatim}

Possible output:

\begin{verbatim}
words: {'eggs', 'spam', 'larch'}
type(words): <class 'set'>
len(words): 3
\end{verbatim}

The source expression contains five textual entries. The set contains only three elements because \texttt{"spam"} is equal to the other \texttt{"spam"} values.

This behavior is not a later cleanup pass. It happens during insertion. Each candidate element is checked against the set's existing contents according to hash-table rules.

A list preserves duplicates because repetition is part of list structure.

\begin{verbatim}
word_list = ["spam", "eggs", "spam", "spam", "larch"]
word_set = {"spam", "eggs", "spam", "spam", "larch"}

print(f"word_list: {word_list!r}")
print(f"len(word_list): {len(word_list)}")
print()

print(f"word_set:  {word_set!r}")
print(f"len(word_set):  {len(word_set)}")
\end{verbatim}

Possible output:

\begin{verbatim}
word_list: ['spam', 'eggs', 'spam', 'spam', 'larch']
len(word_list): 5

word_set:  {'eggs', 'spam', 'larch'}
len(word_set):  3
\end{verbatim}

The list answers: ``which values appear at which positions?''

The set answers: ``which distinct values are members of this domain?''

The same rule applies to numeric values.

\begin{verbatim}
nums = {42, 42, 42, 7, 7, 1}

print(f"nums: {nums!r}")
print(f"type(nums): {type(nums)}")
print(f"len(nums): {len(nums)}")
\end{verbatim}

Possible output:

\begin{verbatim}
nums: {1, 42, 7}
type(nums): <class 'set'>
len(nums): 3
\end{verbatim}

The repeated integer \texttt{42} appears only once. Repetition in the source syntax does not create repetition inside the set.

We can inspect the set object again with \texttt{dir()}.

\begin{verbatim}
nums = {42, 7, 1}

print(f"type(nums): {type(nums)}")
print(f"__contains__ in dir(nums): {'__contains__' in dir(nums)}")
print(f"__iter__     in dir(nums): {'__iter__' in dir(nums)}")
print(f"add          in dir(nums): {'add' in dir(nums)}")
print(f"discard      in dir(nums): {'discard' in dir(nums)}")
print(f"copy         in dir(nums): {'copy' in dir(nums)}")
\end{verbatim}

Possible output:

\begin{verbatim}
type(nums): <class 'set'>
__contains__ in dir(nums): True
__iter__     in dir(nums): True
add          in dir(nums): True
discard      in dir(nums): True
copy         in dir(nums): True
\end{verbatim}

The relevant methods here are:

\begin{itemize}
    \item \texttt{\_\_contains\_\_}: checks membership with \texttt{in};
    \item \texttt{\_\_iter\_\_}: permits traversal over the current elements;
    \item \texttt{add}: inserts one element while preserving uniqueness;
    \item \texttt{discard}: removes an element if present;
    \item \texttt{copy}: creates a shallow copy of the set.
\end{itemize}

The \texttt{add()} method also respects the uniqueness invariant.

\begin{verbatim}
names = set()

names.add("Homer")
names.add("Marge")
names.add("Homer")
names.add("Lisa")
names.add("Homer")

print(f"names: {names!r}")
print(f"len(names): {len(names)}")
\end{verbatim}

Possible output:

\begin{verbatim}
names: {'Lisa', 'Homer', 'Marge'}
len(names): 3
\end{verbatim}

Calling \texttt{add("Homer")} multiple times does not create multiple stored set elements. Once an equal element is already present, the set remains structurally unchanged.

\subsubsection{The Hashability Criterion: Value Equality, Hash Stability, and Table Placement Requirements}

A set element must be hashable. In practical terms, Python must be able to call \texttt{hash()} on the object, and the object's hash-relevant state must remain stable while the object is stored.

Immutable built-in values such as integers, strings, and tuples of hashable elements are accepted.

\begin{verbatim}
item1 = 42
item2 = "D'oh!"
item3 = ("Arthur", "Brian")

print(f"item1: {item1!r}, hash: {hash(item1)}")
print(f"item2: {item2!r}, hash: {hash(item2)}")
print(f"item3: {item3!r}, hash: {hash(item3)}")
\end{verbatim}

Possible output:

\begin{verbatim}
item1: 42, hash: 42
item2: "D'oh!", hash: -5190362190904680204
item3: ('Arthur', 'Brian'), hash: 7652311734482103245
\end{verbatim}

The exact string and tuple hash values may differ between Python executions. The important point is that these objects can produce hash values during the running process.

These values can be inserted into a set.

\begin{verbatim}
items = {42, "D'oh!", ("Arthur", "Brian")}

print(f"items: {items!r}")
print(f"type(items): {type(items)}")
\end{verbatim}

Possible output:

\begin{verbatim}
items: {42, ('Arthur', 'Brian'), "D'oh!"}
type(items): <class 'set'>
\end{verbatim}

Mutable built-in containers such as \texttt{list}, \texttt{dict}, and \texttt{set} are not accepted as set elements.

\begin{verbatim}
bad_item = ["spam", "eggs"]

print(f"bad_item: {bad_item!r}")
print(f"type(bad_item): {type(bad_item)}")
print(f"bad_item.__hash__: {bad_item.__hash__}")

try:
    values = {bad_item}
except TypeError as err:
    print(f"set construction failed: {err}")
\end{verbatim}

Possible output:

\begin{verbatim}
bad_item: ['spam', 'eggs']
type(bad_item): <class 'list'>
bad_item.__hash__: None
set construction failed: unhashable type: 'list'
\end{verbatim}

The method name \texttt{\_\_hash\_\_} exists on the object, but its value is \texttt{None}. This means Python deliberately marks the object as unhashable.

The reason is not that Python is unable to invent an integer for a list. The reason is that a list can change its equality-relevant contents after insertion.

\begin{verbatim}
items = ["spam", "eggs"]

print(f"before: {items!r}")
items.append(42)
print(f"after:  {items!r}")
\end{verbatim}

Possible output:

\begin{verbatim}
before: ['spam', 'eggs']
after:  ['spam', 'eggs', 42]
\end{verbatim}

If the list were used as a set element and its contents changed, the set could lose the ability to locate it correctly. The table slot would have been chosen from one state, while later lookup would describe another state.

A tuple is accepted only if all elements inside it are hashable.

\begin{verbatim}
safe_tuple = ("spam", "eggs", 42)
bad_tuple = ("spam", ["eggs", 42])

print(f"safe_tuple: {safe_tuple!r}")
print(f"hash(safe_tuple): {hash(safe_tuple)}")
print()

print(f"bad_tuple: {bad_tuple!r}")

try:
    print(hash(bad_tuple))
except TypeError as err:
    print(f"hash failed: {err}")
\end{verbatim}

Possible output:

\begin{verbatim}
safe_tuple: ('spam', 'eggs', 42)
hash(safe_tuple): 6766429815130234107

bad_tuple: ('spam', ['eggs', 42])
hash failed: unhashable type: 'list'
\end{verbatim}

The outer tuple is immutable, but it contains a mutable list. Since tuple equality depends on its contained elements, and the contained list is not hashable, the tuple cannot be safely used as a set element.

The ingestion rule can be stated directly:

\begin{quote}
A set element must be usable as a stable hash-table descriptor. It must not change the data that determines its hash and equality while it is stored.
\end{quote}

\subsubsection{CPython Open-Addressing Architecture: Hash Table Slots, Collision Probing, Empty Slots, and Deleted-Entry Markers}

At the implementation level, CPython sets are based on a hash table. The table contains slots. A slot may be empty, may contain an active element, or may hold an internal marker left behind by a deletion.

The exact implementation is CPython-specific and may change between versions, but the useful model is stable enough for understanding ordinary behavior.

A simplified table can be imagined like this:

\begin{verbatim}
slot 0: empty
slot 1: active element
slot 2: active element
slot 3: deleted marker
slot 4: empty
slot 5: active element
\end{verbatim}

When an element is inserted, CPython computes the element's hash and uses part of that hash to choose a starting slot. If that slot is empty, the element can be placed there. If that slot is occupied, CPython probes other slots according to its collision-resolution procedure.

A collision happens when two different elements compete for the same initial table area.

\begin{verbatim}
hash(element_a) -> starting slot 5
hash(element_b) -> starting slot 5

slot 5 already occupied
probe another slot
place element_b in a later acceptable slot
\end{verbatim}

A collision does not mean equality. It only means that the hash-table placement path overlaps. Equality is still checked to decide whether the candidate element is already present.

This is the reason for the two-stage logic:

\begin{verbatim}
first:  use hash to find where to search
second: use equality to confirm the actual match
\end{verbatim}

A simplified membership lookup can be described as follows:

\begin{verbatim}
search_hash = hash(search_value)
start at a table slot derived from search_hash

while the table path continues:
    if the current slot is empty:
        report "not present"
    if the current slot contains an equal element:
        report "present"
    otherwise:
        probe another slot
\end{verbatim}

The empty slot is important. It tells the lookup that the searched element is not present along that probe path. If the search reaches a genuinely empty slot, there is no later element to find for that lookup path.

Deletion complicates this slightly. Suppose an element was removed from the middle of a probe path. CPython cannot always turn that slot into a normal empty slot immediately, because doing so could break future lookups for other elements placed later in the same collision path.

So implementations use a deleted-entry marker, often described as a dummy marker.

\begin{verbatim}
before deletion:
slot 5: element_a
slot 6: element_b
slot 7: empty

after deleting element_a:
slot 5: deleted marker
slot 6: element_b
slot 7: empty
\end{verbatim}

During lookup, the deleted marker does not count as a successful match, but it also does not stop the search in the way a truly empty slot would.

\begin{verbatim}
lookup path:
    active slot      -> compare equality
    deleted marker   -> keep probing
    empty slot       -> stop, not present
\end{verbatim}

This is why removal in a hash table is not the same operation as removing an item from a list. A list deletion shifts later components left to close the positional gap. A set deletion must preserve the correctness of hash-table lookup paths.

The following program does not expose internal slots directly, because Python deliberately hides those implementation details. It demonstrates the external behavior that the internal engine supports.

\begin{verbatim}
names = {"Homer", "Marge", "Lisa", "Bart"}

print(f"start: {names!r}")
print(f"{'Lisa' in names = }")

names.remove("Lisa")

print(f"after remove: {names!r}")
print(f"{'Lisa' in names = }")
print(f"{'Bart' in names = }")
\end{verbatim}

Possible output:

\begin{verbatim}
start: {'Lisa', 'Homer', 'Bart', 'Marge'}
'Lisa' in names = True
after remove: {'Homer', 'Bart', 'Marge'}
'Lisa' in names = False
'Bart' in names = True
\end{verbatim}

After removing \texttt{"Lisa"}, other elements must remain findable. Internally, this is why deletion must respect the table's probing structure.

The important architectural picture is:

\begin{itemize}
    \item hash values guide the starting position;
    \item collisions are resolved by probing;
    \item equality confirms actual element identity by value;
    \item empty slots can terminate failed searches;
    \item deleted-entry markers preserve lookup paths after removal.
\end{itemize}

\subsubsection{Algorithmic Efficiency Matrix: Amortized O(1) Membership Lookups vs. O(N) Sequential Traversal}

The reason sets exist as a separate structure is not only syntax. Their hash-table architecture changes the expected cost of common operations.

The clearest comparison is membership testing.

In a list, membership testing is sequential. Python may need to compare the searched value against many elements.

\begin{verbatim}
items = ["Homer", "Marge", "Lisa", "Bart", "Maggie"]

print(f"{'Maggie' in items = }")
print(f"{'Burns' in items = }")
\end{verbatim}

Conceptually, a failed list membership test looks like this:

\begin{verbatim}
compare with "Homer"
compare with "Marge"
compare with "Lisa"
compare with "Bart"
compare with "Maggie"
not found
\end{verbatim}

If the list has \texttt{N} elements, a failed membership check may require \texttt{N} comparisons. This is called \texttt{O(N)} traversal.

A set uses hash-derived lookup instead.

\begin{verbatim}
items = {"Homer", "Marge", "Lisa", "Bart", "Maggie"}

print(f"{'Maggie' in items = }")
print(f"{'Burns' in items = }")
\end{verbatim}

Conceptually, set membership looks more like this:

\begin{verbatim}
compute hash("Burns")
go to the table region indicated by the hash
probe if needed
stop when the value is found or the path reaches an empty slot
\end{verbatim}

On average, this is treated as amortized \texttt{O(1)} membership testing. That does not mean literally one CPU instruction. It means the expected amount of work does not grow linearly with the number of stored elements under normal table conditions.

The comparison can be summarized as follows:

\begin{center}
\begin{tabular}{lll}
\textbf{Operation} & \textbf{List} & \textbf{Set} \\
\hline
Membership test with \texttt{in} & Sequential scan & Hash-table lookup \\
Typical membership cost & \texttt{O(N)} & Amortized \texttt{O(1)} \\
Duplicate storage & Preserved & Removed automatically \\
Positional indexing & Supported & Not supported \\
Order as structure & Yes & No \\
Element requirement & Any object reference & Hashable object \\
\end{tabular}
\end{center}

A small script can show the semantic difference, without pretending to be a precise benchmark.

\begin{verbatim}
values_list = ["name_0", "name_1", "name_2", "name_3", "name_4"]
values_set = set(values_list)

target = "name_4"

print(f"values_list: {values_list!r}")
print(f"values_set:  {values_set!r}")
print()

print(f"target: {target!r}")
print(f"target in values_list: {target in values_list}")
print(f"target in values_set:  {target in values_set}")
\end{verbatim}

Possible output:

\begin{verbatim}
values_list: ['name_0', 'name_1', 'name_2', 'name_3', 'name_4']
values_set:  {'name_0', 'name_4', 'name_3', 'name_1', 'name_2'}

target: 'name_4'
target in values_list: True
target in values_set:  True
\end{verbatim}

Both containers answer the membership question correctly. The difference is the access route. The list reaches the answer through sequence traversal. The set reaches the answer through hash-table lookup.

For tiny containers, the practical timing difference may be irrelevant. For large membership-heavy workloads, the structural difference becomes important.

\begin{verbatim}
names = [
    "Arthur",
    "Brian",
    "Lancelot",
    "Galahad",
    "Bedevere",
    "Robin",
]

allowed = set(names)

candidate = "Galahad"

print(f"candidate: {candidate!r}")
print(f"candidate in names:   {candidate in names}")
print(f"candidate in allowed: {candidate in allowed}")
\end{verbatim}

Possible output:

\begin{verbatim}
candidate: 'Galahad'
candidate in names:   True
candidate in allowed: True
\end{verbatim}

The list \texttt{names} is appropriate if the original order matters. The set \texttt{allowed} is appropriate if the main question is membership.

The algorithmic lesson is:

\begin{quote}
Use a list when positional sequence is part of the data. Use a set when the program mainly needs a unique membership domain.
\end{quote}

This is not a matter of taste. It follows from the internal access architecture.